\section{Threats to Validity}
\label{sec:threats-validity}

\subsection{Construct validity}
\label{subsec:construct-validity}

A first construct-validity threat is the relationship between the TLA+ and
Alloy representations. Both address the same cross-shard protocol obligations
but use different state, transition, property, and analysis semantics.
Conceptually related properties are therefore not assumed to be syntactically
or semantically identical across tools. Claims are defined at the
protocol-obligation level, with mutation detection evaluated against the
designated property in each encoding. For example, commit after abort maps to
\texttt{TerminalStateIrreversibility} in Alloy and
\texttt{DecisionConsistency} in TLA+.

Verification-cost metrics introduce a second construct threat. Elapsed time and
maximum resident memory are observable for both tools, but state-space measures
are tool-specific. TLC distinct-state counts are not treated as equivalent to
Alloy's internal search representation. Cost trends are therefore interpreted
within each tool, and absolute timing, memory, or state counts are not used to
claim intrinsic tool superiority. Cross-tool observations are used only as
complementary evidence about the same protocol-level verification problem.

RQ3 introduces a third construct threat because it observes a deterministic
trace abstraction rather than the complete Java program state. The measured
construct is conformance of the selected observable abstraction, not semantic
equivalence between implementation and formal specification. The trace format,
scenario catalogue, mapping rules, and replay procedure were frozen before the
definitive campaign. Both valid and deliberately corrupted traces are tested,
and negative cases must match the expected rejection point rather than merely
fail replay.

Finally, mutation sensitivity is an operational proxy for property-suite
adequacy. The ten predefined scientific mutants exercise selected controls in
replay, receipt validation, terminal decisions, quorum enforcement, and timeout
handling, not an exhaustive defect taxonomy. A mutation score of 1.0 therefore
supports sensitivity to the predefined catalogue only and does not establish
complete defect detection or property-suite completeness. This distinction is
important because a perfect score on a targeted catalogue does not estimate
coverage of defect classes that were never represented experimentally.

\subsection{Internal validity}
\label{subsec:internal-validity}

Host load and execution conditions may affect verification time and memory.
The definitive campaign therefore used a dedicated native Linux environment,
with tasks executed serially at maximum parallelism one, reducing interference
between verification processes. This setup reduces, but cannot eliminate,
variation caused by operating-system scheduling and background activity.

Startup and runtime effects are addressed by separating warm-up executions from
measured repetitions. Warm-ups remain in the execution record but are excluded
from the summaries used to characterize measured verification cost.

Scenario generation and execution ordering could introduce nondeterminism.
Deterministic seeds, predefined configuration families, and the definitive
trace seed set were frozen before execution, and scenario definitions were not
changed in response to observed outcomes.

Measurement instrumentation is another possible confounder because
operating-system measurements and formal-tool parsers can fail independently
of the scientific execution. Instrumentation errors are therefore classified
separately, and reruns are permitted only under the predefined instrumentation
failure rule, not because a scientific outcome is unfavorable. This separation
prevents tool or instrumentation failures from being silently interpreted as
scientific results and prevents result-dependent reruns.

Timeout handling also affects internal validity. A run reaching the
1800-second limit has unknown eventual completion time and final verification
outcome. Timeout and out-of-memory results are retained as censored
observations, not converted to threshold-valued completions or successful
verification. The TLC-large timeouts therefore remain visible rather than
being selectively rerun.

More generally, the frozen protocol fixes tool versions, bounds, configuration
families, repetitions, seeds, execution order, resource limits, and treatment
of incomplete executions before interpretation of the definitive results,
reducing opportunities for result-dependent adaptation. It does not eliminate
environmental variation, but it makes the corresponding execution conditions
explicit and reproducible.

\subsection{External validity}
\label{subsec:external-validity}

The study evaluates one cross-shard commit protocol and its executable and
formal representations. Its findings characterize that protocol and
experimental implementation, not the full space of cross-shard mechanisms.
They should not be generalized directly to production blockchain systems with
different consensus protocols, network assumptions, cryptographic mechanisms,
execution environments, or adversarial models. The Java implementation is an
experimental simulator and validation artifact rather than a production
client.

The small, medium, and large formal profiles provide controlled bounded
configurations, not arbitrarily large deployments. Completion without a
counterexample within those profiles does not imply tractability or absence of
counterexamples under substantially larger bounds or different topologies. The
profiles vary shards, concurrent transfers, validators, quorum requirements,
receipt copies, and formal-analysis bounds only within the predefined campaign.

The mutation catalogue is likewise limited. Ten predefined scientific mutants
exercise selected protocol controls, while other defect classes and
interactions remain outside the evaluated catalogue. Mutation results therefore
generalize only to the evaluated defect classes and closely related protocol
obligations, not arbitrary bugs or adversarial behaviors.

Trace-conformance results are also scenario-bounded. The experiment uses ten
valid and ten deliberately corrupted scenario classes, each under thirty
deterministic seeds. This catalogue exercises multiple normal, timeout, replay,
quorum, concurrency, and corruption behaviors but is not an exhaustive sample
of Java executions. RQ3 therefore supports claims only for the evaluated trace
classes under the declared abstraction and bounds.

Finally, verification-cost observations depend on the selected models, tool
versions, bounds, failure configurations, execution environment, and resource
policy. They identify a scalability boundary for the evaluated setup, not
universal complexity laws for TLA+, Alloy, or cross-shard verification.
Accordingly, claims are restricted to the evaluated protocol, bounds,
scenarios, mutants, and tool configurations. The preserved protocol and
artifact support reproduction of those conditions, but do not expand the
population to which the results generalize.

\subsection{Conclusion validity}
\label{subsec:conclusion-validity}

RQ1 reports no counterexample among completed bounded verification executions,
but this is neither a statistical nor an unbounded proof of correctness. The
fully censored TLC-large profile prevents treating all scheduled RQ1 executions
as completed checks without violations. The RQ1 conclusion therefore applies
only to completed bounded configurations and retains censoring explicitly,
rather than treating incomplete executions as successful property checks.

Detection of all ten predefined scientific mutants yields a mutation score of
1.0 for the evaluated catalogue, but the mutants are not a random sample from
the universe of protocol or implementation defects. The score is descriptive
of sensitivity to the injected mutants, not an estimator of arbitrary future
defect-detection probability.

RQ3 likewise reports perfect observed classification for the twenty predefined
trace classes across thirty deterministic seeds per class. These seeds and
scenario classes are not a random sample of all possible Java executions.
Observed proportions and interval estimates therefore characterize the
evaluated experimental population rather than guaranteeing the same behavior
for arbitrary traces. The perfect observed classification rates therefore do
not constitute a population-level guarantee beyond the frozen catalogue.

RQ4 is limited by sparse and censored profile-level data. Alloy provides only
three completed ordered profiles. A Spearman rank association of 1.0 across
three levels establishes the observed monotonic ordering but cannot determine a
functional growth form or establish nonlinear scaling. TLC provides completed
measurements only for the small and medium profiles, while all measured
large-profile executions are censored by timeout, preventing a comparable
three-level analysis or estimation of a large-profile completion-time
distribution.

Censoring also affects descriptive summaries because medians computed from
completed executions characterize only those completions. If more demanding
runs are preferentially censored, completed-run summaries may underestimate the
computational burden of the full configuration. This is especially relevant
when censoring concentrates in the more demanding profiles, as observed for
TLC-large.

H4 is therefore interpreted as evidence of increasing verification cost
without sufficient evidence for the preregistered nonlinear-growth claim. No
asymptotic complexity law is inferred. More generally, conclusions emphasize
bounded property outcomes, mutation sensitivity, trace classification, effect
direction, and transparent censoring rather than significance claims not
supported by the frozen experimental design.
